\section*{РЕФЕРАТ}
\addcontentsline{toc}{section}{РЕФЕРАТ}

Отчёт о научно-исследовательской работе содержит \pageref{LastPage}~с.,
\ref{lastfig}~рис., \ref{lasttab}~табл., \ref{lastref}~источников,
3~приложения.

\textbf{Ключевые слова:} обнаружение вторжений, NIDS, NTA, анализ
сетевого трафика, глубокое обучение, нейронные сети, CNN-LSTM,
UNSW-NB15, CICIDS2017, ИИ-злоумышленник, дрейф распределений, обнаружение
аномалий, фокальная функция потерь, температурное масштабирование, PSI, MMD.

\textbf{Объект исследования} --- сетевой трафик корпоративной сети,
подверженной несанкционированным действиям и аномальной активности.

\textbf{Предмет исследования} --- методика обнаружения аномальной
активности и несанкционированных действий с использованием нейронных
сетей и активного тестирования посредством ИИ-злоумышленника.

\textbf{Цель работы} --- разработка методики обнаружения аномальной
активности и несанкционированных действий в корпоративной сети на
основе нейронных сетей, включающей анализ сетевого трафика,
классификацию угроз и моделирование атак с помощью ИИ-злоумышленника,
генерирующего трафик в реальном времени для тестирования системы.

\textbf{Методы исследования.} Систематический анализ литературы по
NIDS и глубокому обучению; математическое моделирование атак
параметризованными шаблонами с конечным автоматом планирования и
drift-инжектором; обучение нейронных сетей и классических ML-моделей
в равных условиях с фиксированной схемой препроцессинга;
статистический анализ результатов через bootstrap-доверительные
интервалы 95\,\% и парный критерий Уилкоксона.

\textbf{Основные результаты.} Сформирована методика обнаружения,
включающая препроцессинг flow-данных, библиотеку из семи
параметризованных шаблонов атак, конечный автомат планирования режимов
\emph{норма/атака/дрейф} (11 состояний) и блок drift-monitoring на
основе индексов PSI, KL и MMD. Реализован программный прототип на
стеке Python\,/\,PyTorch с интерфейсами FastAPI и Streamlit (более
4200~строк кода, 9 модулей юнит-тестов). На датасете UNSW-NB15
проведено шесть экспериментов (E1--E6) с 22 запусками и 14 моделями.
Лучшее F1 (среднее по трём seed) --- 0{,}919 у XGBoost; среди
DL-моделей лидируют CNN-LSTM, TCN и BiLSTM с F1 в диапазоне
0{,}77\,--\,0{,}84 при FPR порядка 0{,}01\,--\,0{,}03 после
калибровки порога под целевую частоту ложных срабатываний.
Inference-задержка CNN-LSTM на CPU составила 0{,}96~мс/окно, что
удовлетворяет требованиям эксплуатации в составе SOC. Гипотеза
исследования подтверждена в части роли ИИ-злоумышленника и уточнена
по части превосходства DL над классическими baselines на табличных
flow-данных UNSW-NB15.

\textbf{Степень внедрения.} Программный прототип допускает развёртывание
в составе подсистемы NTA-аналитики корпоративного SOC. В отчёте
сформированы инженерные рекомендации по интеграции и калибровке
порогов под целевой эксплуатационный FPR.

\textbf{Область применения.} Системы обнаружения вторжений
корпоративного класса; учебно-исследовательские стенды по сетевой
безопасности; задачи аудита устойчивости NIDS к дрейфу и
вариативности атак.
